\section{Repaired structural results}

The first result is an identification statement. A weak or strict suitability ordering alone does not determine a cardinal acreage vector. Any monotone transformation preserves the ordering but may alter a proportional rule, and neither the transformation nor the ordering supplies an objective, constraints or optimizer selection. This insufficiency is not a numerical failure; it follows from the mismatch between ordinal and cardinal objects.

Existence and tractability follow under conventional conditions. If $X$ is non-empty and compact, per-acre profits are integrable and the risk-feasible set is non-empty, an optimizer exists. The risk-feasible set is convex, and the finite-scenario problem is linear. These conditions preserve the multi-crop stochastic architecture while making its domain explicit.

Operational constraints can force a universal reversal without a risk mechanism. If $s_i>s_j$ and
\begin{equation}
\sup_{x\in X}x_i+\tau_x<\inf_{x\in X}x_j,
\end{equation}
then every feasible allocation reverses the pair. The bound condition $u_i+\tau_x<\ell_j$ is sufficient. Conversely, if an unconstrained expected-profit optimum is risk-feasible, the risk-constrained optimal value is unchanged and the constrained solution set is the intersection of the unconstrained optimal set with the risk-feasible set. One observed slack optimum does not establish equality of solution sets when optima are multiple.

The complete first-order characterization also blocks an acreage-order shortcut. At an optimum $x^*$, there are non-negative multipliers and a subgradient $d\in\partial r_\alpha(x^*)$ such that
\begin{equation}
0=-\mu+\lambda d+\gamma\mathbf 1+\beta c+G^\top\eta-a+b,
\end{equation}
with primal feasibility, dual feasibility and complementary slackness. For crops $i$ and $j$,
\begin{equation}
\mu_i-\mu_j=\lambda(d_i-d_j)+\beta(c_i-c_j)+(G^\top\eta)_i-(G^\top\eta)_j-(a_i-a_j)+(b_i-b_j).
\label{eq:kktpair}
\end{equation}
The common land multiplier cancels, but budget, shared-constraint and bound terms remain. Equation~\ref{eq:kktpair} characterizes marginal optimality; it does not order acreage levels.

Dependence comparative statics require a distributional premise. If fixed marginals within a named family yield $L_{\theta_1}(x)\preceq_{\mathrm{cx}}L_{\theta_2}(x)$ for every $x$ in a declared domain, then loss-CVaR is weakly ordered there. If the domain is all of $X$, risk-feasible sets are nested and optimal expected profit is weakly non-increasing. No crop-specific acreage conclusion follows without additional uniqueness and selection assumptions.

Accordingly, reversal is represented by sets. Let $g^+_{ij}(\theta)$ and $g^-_{ij}(\theta)$ be the maximum and minimum of $x_j-x_i$ over $S(\theta)$. Possible and universal reversal regions are the sets where $g^+_{ij}$ and $g^-_{ij}$ exceed $\tau_x$. Their boundaries may be empty, disconnected, multiple or interval-valued. A unique threshold requires conditions---including continuity, strict monotonicity and an endpoint crossing---that are not general.

Two teacher-Draft themes survive only with restricted meanings. Pseudo-diversification denotes preregistered discordance between ordinary correlation and lower-tail dependence; it implies neither inclusion nor exclusion and has no welfare content. For information, posterior-contingent value is non-negative when ignoring the signal remains feasible, and optimized value is weakly non-decreasing as action sets expand. The difference between informed and uninformed values need not itself be monotone, strictly positive or supermodular. Complete proofs and counterexample boundaries are provided in the Supplementary Information and the canonical theory package.
